\documentclass[11pt]{article}

\usepackage[utf8]{inputenc}
\usepackage[T1]{fontenc}
\usepackage{lmodern}
\usepackage{geometry}
\usepackage{amsmath,amssymb,amsfonts,amsthm,mathtools}
\usepackage{enumitem}
\usepackage{microtype}
\usepackage{hyperref}

\geometry{margin=1in}
\hypersetup{colorlinks=true, linkcolor=black, urlcolor=blue, citecolor=black}
\setlist[enumerate]{topsep=0.35em,itemsep=0.3em}
\setlist[itemize]{topsep=0.3em,itemsep=0.25em}

\newcommand{\Aether}{\AE{}ther}
\newcommand{\Aflow}{\AE{}ther-flow}
\newcommand{\TheoryProgram}{\Aether{} / \Aflow{} framework}
\newcommand{\Lang}{\mathcal{L}_{0}}
\newcommand{\Pack}{\mathcal{P}}
\newcommand{\Enum}{\mathsf{Enum}}
\newcommand{\Cert}{\mathsf{Cert}}
\newcommand{\Out}{\mathsf{Out}}
\newcommand{\Ledger}{\mathsf{Ledger}}
\newcommand{\Status}{\mathsf{Status}}
\newcommand{\Rel}{\mathsf{Rel}}
\newcommand{\GammaB}{\Gamma}
\newcommand{\Inv}{\mathsf{Inv}}

\newtheorem{test}{Benchmark-independence test}
\newtheorem{finding}{Refutation finding}
\newtheorem{variation}{Admissible variation}
\newtheorem{blocker}{Promotion blocker}
\newtheorem{verdict}{Local verdict}
\newtheorem{nextburden}{Next burden}

\title{\textbf{Localization Source-Basis Axiom Selector-Domain Finite Survey Protocol Benchmark-Independence Refutation}\\[0.35em]
\large Invariance Test for Sealed Packets, Enumerators, Certificates, Output Maps, and Status Transitions}
\author{Refuter AgentJob AJ-RT-20260608-021-001}
\date{June 9, 2026}

\begin{document}

\maketitle

\begin{abstract}
This draft tests whether the finite selector-domain survey protocol is
benchmark-independent under admissible variation of the sealed packet
\(\Pack=(R,B,\Omega,\Rel,\GammaB)\), finite enumerator grammars,
completeness certificates, output maps, invariance ledgers, and
\(D_F\) status-transition rules. The result is local and negative. The
protocol is a real control improvement: it makes finite survey machinery
auditable and forbids direct target metric, null-cone, coordinate,
observer-protocol, rank, and proper-time primitives. However, as stated, it
does not prove that the sealed record set, locality bound, source-operation
grammar, relabeling quotient, certificate exclusions, output maps, or status
transitions are source-forced or invariant across benchmark-blind
alternatives. Candidate reconstruction, canonical ontology edit, benchmark
promotion, and Gate Chair review remain blocked.
\end{abstract}

\section{Scope and Refutation Standard}

This artifact is a draft/control output for task \texttt{RT-20260608-021}.
It does not edit canonical ontology TeX, adopt A0--A4 or selection laws as
canonical ontology, authorize candidate reconstruction, issue a Gate Chair
verdict, or reject the broader \TheoryProgram{}.

Let \(\Lang\) denote the registered source language before the A0--A4
extension sequence. Let the finite survey protocol define
\[
\mathcal{S}_i=(\Enum_i(\Pack),\Cert_i,\Out_i,\Ledger_i),
\qquad i\in\{0,1,2,3,4,F\}.
\]
The protocol passes this Refuter test only if localization-relevant outputs
and statuses are invariant under admissible benchmark-blind variation of
\(\Pack\), \(\Enum_i\), \(\Cert_i\), \(\Out_i\), and \(\Ledger_i\), or if
the protocol explicitly preserves all non-invariant families as negative
control outputs.

\begin{test}[Finite-protocol benchmark independence]
For every admissible pair of source-defined finite survey protocols
\(\mathcal{S}\) and \(\mathcal{S}'\) that differ only by benchmark-blind
choices of source records, locality bounds, operation grammars, relabeling
quotients, certificate exclusions, output maps, or status-transition rules,
the localization-relevant output families must agree up to the declared
source relabeling quotient. If they do not agree, the correct local status is
\(\mathsf{family\_valued}\), \(\mathsf{underdetermined}\),
\(\mathsf{benchmark\_sensitive}\), or \(\mathsf{not\_computed}\), not
candidate support.
\end{test}

Sealing a packet proves immutability after commitment. It does not prove that
the committed packet was the unique source-forced packet.

\section{Sealed Packet Variation}

\begin{finding}[A sealed packet is not a uniqueness theorem]
The protocol requires \(\Pack=(R,B,\Omega,\Rel,\GammaB)\) to be committed
before benchmark comparison. That is necessary for reproducibility. It does
not establish that \(R\), \(B\), \(\Omega\), \(\Rel\), or \(\GammaB\) are
forced by the registered source language.
\end{finding}

\begin{variation}[Record and bound variation]
Two sealed packets may use the same source vocabulary while differing in the
finite source-record set \(R\) or locality bound \(B\). One packet can include
only records near a target-compatible neighborhood surrogate; another can
include a broader source-accessible shell or a stricter source support
radius. Both can be sealed before benchmark comparison. Their induced
presentation boundaries, transition fronts, rank spectra, transport cycles,
and status rows need not agree.
\end{variation}

\begin{blocker}[Packet-family invariance not established]
The current protocol does not yet provide a benchmark-blind packet-family
enumerator or a theorem proving invariance over admissible packet families.
Therefore the sealed-input rule blocks mutation but does not establish
benchmark independence.
\end{blocker}

\section{Enumerator Grammar Variation}

\begin{finding}[Finite grammars can underselect source-admissible cases]
The protocol defines finite enumerators for \(D_0\)--\(D_4\). It requires
source-named operations and forbids target primitives. The remaining problem
is underselection: a finite grammar can omit source-definable operations or
quotients that would produce non-target-like outputs.
\end{finding}

\begin{variation}[Operation and quotient variation]
For \(D_0\) and \(D_1\), one grammar may close under immediate source-flow
support while another closes under a two-step source-flow shell. For \(D_2\)
and \(D_3\), one grammar may use exact recurrence words while another uses
prefix or overlap support. For \(D_4\), one grammar may compare cycles only
under exact support while another allows declared-overlap support. These
choices can remain source-defined and benchmark-blind, yet they can change
fronts, first-response ranks, rank spectra, transported cycle ratios, and
holonomy reports.
\end{variation}

\begin{blocker}[Enumerator-family invariance not established]
The finite protocol states the right survey obligation but does not yet
survey a complete family of admissible finite grammars. It also does not
prove that omitted grammars are source-invalid rather than inconvenient for a
target-like result.
\end{blocker}

\section{Certificate Variation}

\begin{finding}[Completeness certificates can encode exclusions]
\(\Cert_i\) is intended to record grammar completeness, duplicate quotients,
and omitted-case classifications. A certificate is evidence only if its
omissions are reproducible from source criteria. Otherwise, it can convert
non-invariant source cases into absent cases.
\end{finding}

\begin{variation}[Omitted-case classification variation]
One certificate can classify a support relation as duplicate under
\(\Rel\). Another can keep it as an incomparable source-admissible case
because it changes a localization-relevant output. Both certificates can be
formally finite and internally consistent. The difference changes whether the
ledger reports an invariant output or a family-valued output.
\end{variation}

\begin{blocker}[Certificate-exclusion invariance not established]
The protocol does not yet contain a certificate-audit theorem requiring every
duplicate, out-of-bound, or syntactically inadmissible exclusion to be stable
under admissible source-defined certificate alternatives.
\end{blocker}

\section{Output-Map Variation}

\begin{finding}[Localization-relevant outputs are still selected]
\(\Out_i\) records boundaries, fronts, ranks, rank spectra, transported
cycle ratios, holonomy reports, and status evidence. The protocol requires
these maps to be source-defined. It does not prove that the choice of which
source quantities are localization-relevant is source-forced rather than
chosen because those quantities resemble the target benchmark.
\end{finding}

\begin{variation}[Observable-basis variation]
One output map may record only boundary and first-response-rank data. Another
may also record adverse transport-cycle ratios or non-singleton quotient
spectra. Both maps can cite source records. If the narrower map yields an
apparent invariant and the broader map yields a family-valued result, the
protocol is not benchmark-independent.
\end{variation}

\begin{blocker}[Output-map stability not established]
Candidate support cannot be drawn from the current protocol until the
output-map basis is either source-forced or shown stable over admissible
benchmark-blind output bases. A favorable projection of a non-invariant
family is benchmark-sensitive.
\end{blocker}

\section{Status-Transition Variation}

\begin{finding}[Status rows improve provenance but do not prove status]
\(D_F\) records packet identifiers, grammar hashes, certificate status,
output-family hashes, target exclusions, comparison time, reviewer, and final
status. These fields make status claims inspectable. They do not prove that
\(\mathsf{invariant}\) is the correct status.
\end{finding}

\begin{variation}[Transition-rule variation]
One status rule may allow \(\mathsf{family\_valued}\) to become
\(\mathsf{invariant}\) after a source-side irrelevance theorem. Another may
require a stronger proof that every output-map basis yields the same
localization-relevant family. Both rules can be source-side. They can produce
different status ledgers for the same enumerated records.
\end{variation}

\begin{blocker}[Status-transition evidence not established]
The current protocol does not define enough independent evidence gates for
status promotion from \(\mathsf{not\_computed}\) or
\(\mathsf{family\_valued}\) to \(\mathsf{invariant}\). Status transitions
therefore remain claim-bearing steps, not neutral metadata updates.
\end{blocker}

\section{Independence Ledger}

\begin{center}
\begin{tabular}{p{0.23\linewidth}p{0.40\linewidth}p{0.26\linewidth}}
\textbf{Surface} & \textbf{Unresolved admissible variation} & \textbf{Local status} \\
\hline
\(\Pack\) & record set, locality bound, operation grammar, quotient, blinding ledger & not benchmark-independent \\
\(\Enum_i\) & finite grammar breadth and source-operation closure & not benchmark-independent \\
\(\Cert_i\) & duplicate, out-of-bound, and inadmissible-case classifications & not benchmark-independent \\
\(\Out_i\) & localization-relevant observable basis and adverse-output retention & not benchmark-independent \\
\(\Ledger_i\) & invariance, family-valued, and benchmark-sensitive classification & not benchmark-independent \\
\(D_F\) & status-transition evidence gates and reviewer rules & blocker metadata incomplete \\
\end{tabular}
\end{center}

\section{Local Verdict}

\begin{verdict}[Finite survey protocol fails the current benchmark-independence test]
The finite survey protocol is an important audit improvement, but it does
not yet establish benchmark independence. It fixes a reproducible machinery
for one sealed survey path; it does not prove that the machinery is invariant
over admissible benchmark-blind alternatives. Hidden target import and
benchmark sensitivity can still enter through packet-family choice,
enumerator-family underselection, certificate exclusions, output-map
projection, and status-transition evidence.
\end{verdict}

This is a local negative result. It does not reject the \TheoryProgram{} and
does not reject finite selector-domain surveys as research objects. It blocks
candidate reconstruction from the present protocol.

\section{Next Burden}

\begin{nextburden}[Finite-variation repair packet or controlled pause]
The logical next step is not Candidate Constructor and not Gate Chair review.
The next bounded task should route to an Ontology Formalizer to propose a
finite-variation repair packet that includes:
\begin{enumerate}
    \item a benchmark-blind family generator for sealed packets
    \((R,B,\Omega,\Rel,\GammaB)\);
    \item enumerator-family completeness obligations for \(D_0\)--\(D_4\);
    \item certificate-exclusion proof obligations;
    \item output-map stability criteria that preserve adverse outputs; and
    \item \(D_F\) evidence gates for every status transition.
\end{enumerate}
If those burdens require authority beyond finite source-defined survey
machinery, the Director should enter a controlled pause rather than
authorize candidate reconstruction.
\end{nextburden}

\section*{References}

Ricciardi, A. S. (2026, June 9). \emph{Localization source-basis axiom
selector-domain finite survey protocol smuggling audit: Target-import review
of sealed inputs, enumerators, certificates, and ledgers} [Unpublished
manuscript]. The Aether-Flow Interpretation of Relativity Research Project,
\texttt{research\_control/tasks/RT-20260608-020/artifacts/19\_LOCALIZATION\_SOURCE\_BASIS\_AXIOM\_SELECTOR\_DOMAIN\_FINITE\_SURVEY\_PROTOCOL\_SMUGGLING\_AUDIT.tex}.

Ricciardi, A. S. (2026, June 9). \emph{Localization source-basis axiom
selector-domain finite survey protocol packet: Source-fixed enumerators and
reproducible status ledgers for D0--DF} [Unpublished manuscript]. The
Aether-Flow Interpretation of Relativity Research Project,
\texttt{research\_control/tasks/RT-20260608-019/artifacts/18\_LOCALIZATION\_SOURCE\_BASIS\_AXIOM\_SELECTOR\_DOMAIN\_FINITE\_SURVEY\_PROTOCOL\_PACKET.tex}.

Ricciardi, A. S. (2026, June 9). \emph{Localization source-basis axiom
selector-domain benchmark-independence refutation: Invariance test for
repaired survey domains D0--DF} [Unpublished manuscript]. The Aether-Flow
Interpretation of Relativity Research Project,
\texttt{research\_control/tasks/RT-20260608-018/artifacts/17\_LOCALIZATION\_SOURCE\_BASIS\_AXIOM\_SELECTOR\_DOMAIN\_BENCHMARK\_INDEPENDENCE\_REFUTATION.tex}.

\end{document}
